\scanpage{25}
\uline{1.1} Soit $n\in\mathbb N^*$, et posons
$C_n\subset\mathbb P^2_{\mathbb Q}$ équation projective
\[
x^n+y^n+z^n=0\tag*{(1)}
\]
Le groupe $\mu_{n\,\mathbb Q}^{3}$ opère sur $\mathbb E^3_{\mathbb Q}$ par matrices diagonales,
donc posant
\[
\unsure{\pi|\Pi}_{n\,\mathbb Q}=\mu_n^3/\underset{\text{\scriptsize$\hookrightarrow$ diag.}}{\mu_n}\simeq
\mu_n\otimes_{\mathbb Z}(\mathbb Z^3/\mathbb Z)\tag*{(2)}
\]
\margin{NB on écrit $\mu_n$ au lieu de $\mu_{n\,\mathbb Q}$}
opère sur $\mathbb P^2_{\mathbb Q}$, en laissant $C_n$ invariant. La
courbe $C_n/\pi_{n\,\mathbb Q}$ s'identifie à la courbe \inserted{$C_1$} qui est
donnée, en coordonnées homogènes, par
\[
X+Y+Z=0\tag*{(3)}
\]
qu'on identifie à $\mathbb P^1_{X}$, en envoyant les
points de $C_1$ pour lesquels $X,Y$ ou $Z$ sont
nuls resp. en $0,1,\infty$. Donc écrivant $C_1$
en coordonnées affines, faisant $Z=1$
i.e. $X/Z=\deleted{x}\xi$, $Y/Z=\deleted{y}\eta$, sous la forme
\[
\deleted{x}\xi+\deleted{y}\eta+1=0\tag*{($3'$)}
\]
i.e.
\[
\deleted{y}\eta=-1-\deleted{x}\xi
\]
$\deleted{x}\xi=0$ correspond à $\deleted{y}\eta=-1$, $\deleted{y}\eta=0$ à $\deleted{x}\xi=-1$, donc
posant
\[
t=-\deleted{x}\xi\quad\text{donc}\quad \xi=-t,\ \deleted{y}\eta=t-1,\ \zeta=1\tag*{($3''$)}
\]
les valeurs de $t$ pour lesquels $X=0,Y=0,Z=0$
sont resp. $t=0,1,\infty$.

$C_n$ est étale sur $C_1=\mathbb P^1_{\mathbb Q}$ sauf en $0,1,\infty$, on
trouve rev. \deleted{\unsure{pr}} principal $C_n^*$ de $C_1^*\simeq\mathbb P^1_{\mathbb Q}\setminus\{0,1,\infty\}$, groupe $\pi_{n\,\mathbb Q}$.
% margin note written sideways
\margin{Il y a en $\alpha(i)$ $(i=0,1,\infty)$ $n$ pts de ramification de $C_n$, chacun de degré de ramification $n$.}
% margin note written sideways
\margin{Comme carte pondérée, elle est formée de $2n^2$ triangles chaque sommet est d'ordre $2n$, chaque type de sommet \unsure{intervient} \unsure{$n$|en} fois\ldots}
% archive sheet number: 24
